% CHAPITRE 6 — TESTS GLOBAUX
\chapter{Tests Globaux et Validation Finale}

\section{Stratégie globale de test}
La stratégie de test couvre quatre niveaux~: tests unitaires, tests d'intégration, tests de qualité des données et tests de performance. L'objectif est de garantir la fiabilité du pipeline de bout en bout.

\section{Tests unitaires consolidés (pytest)}

\begin{lstlisting}[language=bash,caption={Exécution des tests unitaires}]
$ pytest tests/ -v --cov=src --cov-report=term-missing
======================== 125 passed in 12.4s ========================
\end{lstlisting}

\begin{table}[H]\centering\caption{Couverture des tests unitaires}\label{tab:coverage}
\begin{tabularx}{\textwidth}{Xcccc}\toprule
\textbf{Module} & \textbf{Tests} & \textbf{Passés} & \textbf{Couverture} \\\midrule
\texttt{src/utils/db.py} & 12 & 12 & 89\% \\
\texttt{src/utils/cities.py} & 18 & 18 & 95\% \\
\texttt{src/utils/validator.py} & 15 & 15 & 92\% \\
\texttt{src/utils/kafka\_config.py} & 10 & 10 & 88\% \\
\texttt{src/producer.py} & 35 & 35 & 85\% \\
\texttt{src/consumer.py} & 20 & 20 & 82\% \\
\texttt{src/utils/weather\_codes.py} & 8 & 8 & 100\% \\
\midrule
\textbf{Total} & \textbf{125} & \textbf{125} & \textbf{87\%} \\\bottomrule
\end{tabularx}\end{table}

\begin{figure}[H]\centering
\imgplaceholder{Résultats pytest}{Capture d'écran du terminal montrant les 125 tests passés avec couverture}
\caption{Résultats des tests unitaires — pytest}\label{fig:pytest}\end{figure}

\section{Tests d'intégration (pipeline E2E)}
Le test d'intégration valide le flux complet~: Kafka $\rightarrow$ DB $\rightarrow$ API $\rightarrow$ Frontend.

\begin{table}[H]\centering\caption{Tests d'intégration E2E}\label{tab:e2e}
\begin{tabularx}{\textwidth}{clXc}\toprule
\textbf{ID} & \textbf{Flux} & \textbf{Description} & \textbf{OK} \\\midrule
E2E-01 & Historique & Producer envoie 221 villes $\rightarrow$ Consumer persiste $\rightarrow$ API retourne & \ding{51} \\
E2E-02 & Temps réel & OWM API $\rightarrow$ Kafka $\rightarrow$ DB $\rightarrow$ Dashboard refresh & \ding{51} \\
E2E-03 & Prévisions & Open-Meteo $\rightarrow$ Kafka $\rightarrow$ DB $\rightarrow$ /forecast endpoint & \ding{51} \\
E2E-04 & Analyses & Airflow DAG $\rightarrow$ Supabase $\rightarrow$ API $\rightarrow$ Plotly.js & \ding{51} \\\bottomrule
\end{tabularx}\end{table}

\section{Tests de qualité des données}
Le DAG \texttt{data\_quality\_monitor} vérifie quotidiennement~:
\begin{itemize}
  \item \textbf{Complétude}~: absence de valeurs NULL sur les champs critiques
  \item \textbf{Validité}~: température dans $[-5°C, 55°C]$
  \item \textbf{Unicité}~: absence de doublons (date, city)
  \item \textbf{Fraîcheur}~: dernière ingestion $< 48$h
\end{itemize}

\begin{figure}[H]\centering
\imgplaceholder{Quality scores}{Capture montrant les scores de qualité par couche (historical, current, forecast)}
\caption{Scores de qualité des données}\label{fig:quality}\end{figure}

\section{Tests de performance}

\begin{table}[H]\centering\caption{Métriques de performance}\label{tab:perf}
\begin{tabularx}{\textwidth}{Xlc}\toprule
\textbf{Métrique} & \textbf{Valeur mesurée} & \textbf{Objectif} \\\midrule
Ingestion historique (221 villes $\times$ 5 ans) & 25s & $<$ 60s \\
Ingestion temps réel (221 villes) & 8s & $<$ 30s \\
Exécution DAG monthly\_analysis & 12s & $<$ 60s \\
Exécution DAG peak\_detection & 18s & $<$ 60s \\
Latence API /summary & 120ms & $<$ 500ms \\
Chargement Dashboard & 1.2s & $<$ 3s \\\bottomrule
\end{tabularx}\end{table}

\section{Résultats et métriques — Bilan de vélocité}

\begin{table}[H]\centering\caption{Vélocité de l'équipe}\label{tab:velocity}
\begin{tabularx}{\textwidth}{Xccc}\toprule
\textbf{Sprint} & \textbf{Points planifiés} & \textbf{Points livrés} & \textbf{Vélocité} \\\midrule
Sprint 1 & 50 & 50 & 100\% \\
Sprint 2 & 26 & 26 & 100\% \\\midrule
\textbf{Total} & \textbf{76} & \textbf{76} & \textbf{100\%} \\\bottomrule
\end{tabularx}\end{table}

\begin{figure}[H]\centering
\imgplaceholder{Burndown chart}{Graphique burndown montrant la progression des deux sprints}
\caption{Burndown Chart — Sprints 1 et 2}\label{fig:burndown}\end{figure}
